\documentclass[aspectratio=169,11pt]{beamer}
\usepackage[utf8]{inputenc}
\usepackage[T1]{fontenc}
\usepackage[spanish,mexico]{babel}
\usepackage{amsmath,amssymb,mathtools}
\usetheme{Madrid}
\usecolortheme{default}
\setbeamertemplate{navigation symbols}{}
\setbeamertemplate{footline}[frame number]
\setbeamerfont{frametitle}{size=\large}
\setbeamersize{text margin left=9mm,text margin right=9mm}
\newcommand{\R}{\mathbb R}
\newcommand{\M}{\operatorname{M}}
\newcommand{\rang}{\operatorname{rang}}
\newcommand{\nul}{\operatorname{nul}}
\title{Álgebra Lineal}
\subtitle{Unidad 3: Matrices y determinantes\\Versión de estudio para teléfono}
\author{Carlos Ernesto Martínez}
\institute{Tecnológico Nacional de México · ACF-0903}
\date{Octubre de 2026}
\begin{document}
\begin{frame}\titlepage\end{frame}
\begin{frame}{Ruta de estudio}
\tableofcontents
\end{frame}
\section{Matrices y clasificación}
\begin{frame}{¿Por qué introducir matrices?}
Un sistema lineal organiza coeficientes en renglones y columnas.
\[\begin{cases}2x-y+3z=5,\\x+4y-z=2\end{cases}\quad\longleftrightarrow\quad A=\begin{pmatrix}2&-1&3\\1&4&-1\end{pmatrix}\]
\begin{itemize}\item $A$ tiene \textbf{2 renglones y 3 columnas}: es de orden $2\times3$.\item No es cuadrada: en general, las matrices pueden ser rectangulares.\item Cada entrada $a_{ij}$ está en el renglón $i$, columna $j$.\end{itemize}
\end{frame}
\begin{frame}{Definición de matriz}
Una matriz $m\times n$ sobre $\R$ es un arreglo rectangular de $mn$ números:
\[A=(a_{ij})_{m\times n}=\begin{pmatrix}a_{11}&\cdots&a_{1n}\\\vdots&\ddots&\vdots\\a_{m1}&\cdots&a_{mn}\end{pmatrix}.\]
\textbf{Ejemplo:} $B=\begin{pmatrix}1&0&4&-2\\3&5&1&0\\-1&2&2&7\end{pmatrix}$ es $3\times4$, con $b_{24}=0$ y $b_{31}=-1$.
\end{frame}
\begin{frame}{Igualdad y tipos}
Dos matrices son iguales sólo cuando \textbf{tienen las mismas dimensiones} y coinciden entrada por entrada.
\[\begin{pmatrix}x&2\\3&y\end{pmatrix}=\begin{pmatrix}5&2\\3&-1\end{pmatrix}\implies x=5,\quad y=-1.\]
Ejemplos de matrices especiales: nula, diagonal, identidad, triangular superior e inferior, simétrica ($A^T=A$) y antisimétrica ($A^T=-A$).
\end{frame}
\begin{frame}{Clasificación: ejemplos distintos}
\[\underbrace{\begin{pmatrix}2&0&0\\0&-1&0\\0&0&3\end{pmatrix}}_{\text{diagonal}}\qquad\underbrace{\begin{pmatrix}1&2&3\\0&4&5\\0&0&6\end{pmatrix}}_{\text{triangular superior}}\]
\[\underbrace{\begin{pmatrix}1&2&-1\\2&0&3\\-1&3&4\end{pmatrix}}_{\text{simétrica}}\]
\textbf{Cuidado:} una matriz rectangular no puede ser simétrica en el sentido habitual.
\end{frame}
\section{Operaciones matriciales}
\begin{frame}{Suma y producto por un escalar}
Para $A,B\in\M_{m\times n}(\R)$:
\[(A+B)_{ij}=a_{ij}+b_{ij},\qquad(\lambda A)_{ij}=\lambda a_{ij}.\]
\textbf{Ejemplo con matrices $2\times3$:}
\[A=\begin{pmatrix}1&-2&3\\0&4&1\end{pmatrix},\quad B=\begin{pmatrix}2&1&0\\-1&3&5\end{pmatrix}\]
\[A+B=\begin{pmatrix}3&-1&3\\-1&7&6\end{pmatrix},\quad2A=\begin{pmatrix}2&-4&6\\0&8&2\end{pmatrix}.\]
\end{frame}
\begin{frame}{Propiedades: ¿de dónde salen?}
Para matrices del mismo orden y escalares $\lambda,\mu$:
\[A+B=B+A,\quad A+(B+C)=(A+B)+C\]
\[\lambda(A+B)=\lambda A+\lambda B,\quad(\lambda+\mu)A=\lambda A+\mu A.\]
\textbf{Demostración breve de conmutatividad:}
\[(A+B)_{ij}=a_{ij}+b_{ij}=b_{ij}+a_{ij}=(B+A)_{ij}.\]
Como coincide cada entrada, las matrices son iguales.
\end{frame}
\begin{frame}{Cuándo existe el producto $AB$}
Si $A$ es $m\times\color{blue}n$ y $B$ es $\color{blue}n\times p$, el producto $AB$ existe y es $m\times p$:
\[(AB)_{ij}=\sum_{k=1}^{n}a_{ik}b_{kj}.\]
\textbf{Idea:} renglón $i$ de $A$ por columna $j$ de $B$.
\[\underbrace{(2\times4)}_A\ \underbrace{(4\times3)}_B\;\longrightarrow\;\underbrace{(2\times3)}_{AB}.\]
Un $3\times4$ multiplicado por un $2\times3$ \textbf{no está definido}.
\end{frame}
\begin{frame}{Producto rectangular, paso 1}
\[A=\begin{pmatrix}1&2&-1\\0&3&2\end{pmatrix},\quad B=\begin{pmatrix}2&0\\1&-1\\3&4\end{pmatrix}\]
$A$ es $2\times3$ y $B$ es $3\times2$; el producto es $2\times2$.
\[(AB)_{11}=1(2)+2(1)+(-1)(3)=1\]
\[(AB)_{12}=1(0)+2(-1)+(-1)(4)=-6.\]
\end{frame}
\begin{frame}{Producto rectangular, paso 2}
Los otros elementos son:
\[(AB)_{21}=0(2)+3(1)+2(3)=9\]
\[(AB)_{22}=0(0)+3(-1)+2(4)=5.\]
Por tanto,
\[\boxed{AB=\begin{pmatrix}1&-6\\9&5\end{pmatrix}}\]
\textbf{Verificación:} todos los elementos son productos escalares de longitud 3.
\end{frame}
\begin{frame}{El producto matricial no conmuta}
En general $AB\neq BA$, incluso si ambos productos existen.
\[A=\begin{pmatrix}1&1\\0&1\end{pmatrix},\quad B=\begin{pmatrix}1&0\\1&1\end{pmatrix}\]
\[AB=\begin{pmatrix}2&1\\1&1\end{pmatrix},\qquad BA=\begin{pmatrix}1&1\\1&2\end{pmatrix}.\]
Sí se conservan la asociatividad y la distributividad cuando las dimensiones son compatibles.
\end{frame}
\begin{frame}{Transpuesta y sus propiedades}
Intercambiamos renglones y columnas: $(A^T)_{ij}=a_{ji}$.
\[A=\begin{pmatrix}1&2&3\\4&5&6\end{pmatrix}\quad\Rightarrow\quad A^T=\begin{pmatrix}1&4\\2&5\\3&6\end{pmatrix}.\]
\[(A+B)^T=A^T+B^T,\quad (AB)^T=B^TA^T.\]
\textbf{Atención:} el orden de los factores se invierte.
\end{frame}
\section{Inversa y Gauss--Jordan}
\begin{frame}{Qué significa tener inversa}
Una matriz cuadrada $A$ es invertible si existe $A^{-1}$ tal que
\[AA^{-1}=A^{-1}A=I.\]
\textbf{No todas las matrices cuadradas son invertibles.}
\[A=\begin{pmatrix}1&2\\2&4\end{pmatrix}\]
La segunda fila es el doble de la primera: $A$ no tiene inversa.
\end{frame}
\begin{frame}{Gauss--Jordan: método general}
Para encontrar $A^{-1}$ formamos la matriz aumentada
\[\left[\begin{array}{c|c}A&I\end{array}\right].\]
Aplicamos operaciones elementales de renglón hasta obtener
\[\left[\begin{array}{c|c}I&A^{-1}\end{array}\right].\]
Si es imposible convertir el bloque izquierdo en $I$, $A$ no es invertible.
\end{frame}
\begin{frame}{Ejemplo de inversa: paso 1}
Sea $A=\begin{pmatrix}2&1\\1&1\end{pmatrix}$.
\[\left[\begin{array}{cc|cc}2&1&1&0\\1&1&0&1\end{array}\right]\xrightarrow{R_1\leftrightarrow R_2}\left[\begin{array}{cc|cc}1&1&0&1\\2&1&1&0\end{array}\right].\]
Ahora hacemos $R_2\leftarrow R_2-2R_1$.
\[\left[\begin{array}{cc|cc}1&1&0&1\\0&-1&1&-2\end{array}\right].\]
\end{frame}
\begin{frame}{Ejemplo de inversa: paso 2}
Multiplicamos $R_2$ por $-1$:
\[\left[\begin{array}{cc|cc}1&1&0&1\\0&1&-1&2\end{array}\right].\]
Hacemos $R_1\leftarrow R_1-R_2$:
\[\left[\begin{array}{cc|cc}1&0&1&-1\\0&1&-1&2\end{array}\right].\]
\[\boxed{A^{-1}=\begin{pmatrix}1&-1\\-1&2\end{pmatrix}}\]
Compruébalo calculando $AA^{-1}=I_2$.
\end{frame}
\section{Determinantes}
\begin{frame}{Determinante de orden dos}
\[\det\begin{pmatrix}a&b\\c&d\end{pmatrix}=ad-bc.\]
\textbf{Ejemplo:}
\[\det\begin{pmatrix}3&-2\\4&1\end{pmatrix}=3(1)-(-2)(4)=11.\]
El determinante es un \textbf{número}, no una matriz. Si vale cero, la matriz cuadrada no es invertible.
\end{frame}
\begin{frame}{Determinante de orden tres}
Para $A=\begin{pmatrix}a&b&c\\d&e&f\\g&h&i\end{pmatrix}$:
\[\det A=aei+bfg+cdh-ceg-bdi-afh.\]
Se puede calcular por Sarrus en orden tres o por cofactores.
\textbf{Advertencia:} la regla de Sarrus \textbf{no} se extiende directamente a matrices $4\times4$.
\end{frame}
\begin{frame}{Menores y cofactores}
El menor $M_{ij}$ se obtiene eliminando el renglón $i$ y la columna $j$ de $A$ y calculando el determinante restante.
\[C_{ij}=(-1)^{i+j}M_{ij}.\]
Los signos alternan:
\[\begin{pmatrix}+&-&+\\-&+&-\\+&-&+\end{pmatrix}.\]
\textbf{Ejemplo:} si $M_{12}=5$, entonces $C_{12}=-5$.
\end{frame}
\begin{frame}{Expansión por cofactores}
A lo largo de cualquier renglón $i$:
\[\det A=\sum_{j=1}^n a_{ij}C_{ij}.\]
\textbf{Consejo práctico:} elige una fila o columna que tenga varios ceros.
\[A=\begin{pmatrix}2&0&0\\3&1&4\\-1&2&5\end{pmatrix}\]
Por la primera fila:
\[\det A=2\det\begin{pmatrix}1&4\\2&5\end{pmatrix}=2(5-8)=-6.\]
\end{frame}
\begin{frame}{Operaciones elementales y determinante}
\begin{itemize}\item Intercambiar dos renglones cambia el signo.\item Multiplicar un renglón por $k$ multiplica el determinante por $k$.\item Sumar un múltiplo de otro renglón no modifica el determinante.\item Para una matriz triangular, el determinante es el producto de las entradas diagonales.\end{itemize}
Estas reglas permiten calcular determinantes de orden 4 o 5 sin expandir decenas de cofactores.
\end{frame}
\begin{frame}{Ejemplo triangular $5\times5$}
\scriptsize
\[A=\begin{pmatrix}2&1&0&3&-1\\0&-1&4&0&2\\0&0&3&1&1\\0&0&0&2&-2\\0&0&0&0&5\end{pmatrix}.\]
\normalsize
Como $A$ es triangular superior,
\[\det A=2(-1)(3)(2)(5)=\boxed{-60}.\]
En cambio, el método de Sarrus no sería aplicable aquí.
\end{frame}
\begin{frame}{Cuándo es invertible una matriz}
Para $A\in\M_{n\times n}(\R)$, son equivalentes:
\begin{itemize}\item $A$ es invertible.\item $\det A\neq0$.\item $\rang A=n$.\item $Ax=0$ sólo tiene la solución $x=0$.\item Las columnas de $A$ son linealmente independientes.\end{itemize}
Esta relación conecta las unidades 2, 3 y 4.
\end{frame}
\section{Adjunta, rango y núcleo}
\begin{frame}{Adjunta e inversa por cofactores}
Formamos la matriz de cofactores $C=(C_{ij})$ y después la transponemos:
\[\operatorname{adj}(A)=C^T.\]
Se cumple
\[A\operatorname{adj}(A)=\det(A)I.\]
Si $\det(A)\neq0$:
\[\boxed{A^{-1}=\dfrac{\operatorname{adj}(A)}{\det A}}.\]
\end{frame}
\begin{frame}{Rango de una matriz rectangular}
El rango es el número de pivotes de una forma escalonada de $A$.
\[A=\begin{pmatrix}1&2&0&1\\0&1&3&2\\0&0&0&0\end{pmatrix}\]
Hay dos pivotes, así que $\rang A=2$. Para una matriz $m\times n$:
\[0\leq\rang A\leq\min\{m,n\}.\]
No hace falta que $A$ sea cuadrada para hablar de rango.
\end{frame}
\begin{frame}{Núcleo: planteamiento}
El núcleo de una matriz $m\times n$ es
\[\ker A=\{x\in\R^n:Ax=0\}.\]
Si $A=\begin{pmatrix}1&2&0\\0&1&3\end{pmatrix}$, entonces
\[x_1+2x_2=0,\qquad x_2+3x_3=0.\]
Tomamos la variable libre $x_3=t$ y resolvemos por sustitución.
\end{frame}
\begin{frame}{Núcleo: solución paso a paso}
De $x_3=t$ y $x_2+3x_3=0$ resulta $x_2=-3t$.
\[x_1+2(-3t)=0\implies x_1=6t.\]
Así,
\[x=\begin{pmatrix}6t\\-3t\\t\end{pmatrix}=t\begin{pmatrix}6\\-3\\1\end{pmatrix},\quad t\in\R.\]
\[\boxed{\ker A=\operatorname{span}\{(6,-3,1)^T\}}.\]
Rango 2, nulidad 1 y $2+1=3$ (número de columnas).
\end{frame}
\begin{frame}{Teorema rango--nulidad}
Para una matriz $A$ de orden $m\times n$,
\[\boxed{\rang A+\dim(\ker A)=n}.\]
La suma siempre da el \textbf{número de columnas}, no el número de renglones.
\textbf{Ejemplo:} una matriz $3\times5$ de rango $3$ tiene nulidad $5-3=2$.
\end{frame}
\section{Práctica organizada}
\begin{frame}{Práctica 1: suma, resta y escalares}
Resuelve individualmente. Comprueba primero las dimensiones.
\small
\begin{enumerate}\item $A=\begin{pmatrix}1&2&0\\-1&3&4\end{pmatrix}$, $B=\begin{pmatrix}2&0&5\\1&-2&3\end{pmatrix}$: calcula $A+B$ y $2A-B$.\item $C=\begin{pmatrix}0&1\\2&-1\\3&4\end{pmatrix}$, $D=\begin{pmatrix}2&3\\1&0\\-2&5\end{pmatrix}$: calcula $3C-D$.\item Si $E$ es $4\times5$ y $F$ es $5\times4$, explica si $E+F$ está definido.\end{enumerate}
\end{frame}
\begin{frame}{Práctica 2: productos rectangulares}
\small
\begin{enumerate}\item $\begin{pmatrix}1&2&0\\3&-1&1\end{pmatrix}\begin{pmatrix}2&1\\0&4\\1&-2\end{pmatrix}$ ($2\times3$ por $3\times2$).\item $\begin{pmatrix}1&0\\2&3\\-1&1\end{pmatrix}\begin{pmatrix}4&2&-1&0\\1&3&2&5\end{pmatrix}$ ($3\times2$ por $2\times4$).\item Decide los órdenes de $AB$ y $BA$ para $A$ de $5\times3$ y $B$ de $3\times5$.\end{enumerate}
\end{frame}
\begin{frame}{Práctica 3: inversas}
Calcula las inversas mediante Gauss--Jordan, o justifica que no existen.
\[A=\begin{pmatrix}1&2\\1&3\end{pmatrix},\quad B=\begin{pmatrix}2&1\\4&2\end{pmatrix},\quad C=\begin{pmatrix}1&0&1\\0&1&2\\0&0&1\end{pmatrix}.\]
Verifica tus resultados multiplicando por la matriz original.
\end{frame}
\begin{frame}{Práctica 4: determinantes}
\small
\begin{enumerate}\item $\det\begin{pmatrix}4&-1\\3&2\end{pmatrix}$.\item $\det\begin{pmatrix}1&2&0\\0&3&1\\2&0&1\end{pmatrix}$ (por cofactores).\item $\det\begin{pmatrix}1&2&0&1\\0&2&3&1\\0&0&-1&2\\0&0&0&3\end{pmatrix}$ (aprovecha su estructura).\end{enumerate}
\end{frame}
\begin{frame}{Práctica 5: rango y núcleo}
Determina rango, una base del núcleo y nulidad.
\[A=\begin{pmatrix}1&2&3\\2&4&6\end{pmatrix},\quad B=\begin{pmatrix}1&0&2&1\\0&1&-1&3\end{pmatrix}.\]
\textbf{Comprobación:} verifica $Ax=0$ con cada vector base y la identidad rango--nulidad.
\end{frame}
\begin{frame}{Dónde está el desarrollo completo}
Esta presentación es una \textbf{guía visual de repaso} para el teléfono.
\begin{itemize}\item Notas autocontenidas: PDF de 72 páginas.\item Problemario: 228 ejercicios agrupados por técnica en el documento completo.\item Respuestas: archivo independiente para autocorrección.\end{itemize}
No reemplaza las demostraciones y los ejemplos extensos de las notas maestras.
\end{frame}
\begin{frame}{Cierre de la unidad}
\centering\Large
Matrices $\longrightarrow$ operaciones $\longrightarrow$ inversas\\[1.4em]
Determinantes $\longrightarrow$ cofactores $\longrightarrow$ adjunta\\[1.4em]
Rango $\longrightarrow$ núcleo $\longrightarrow$ espacios vectoriales
\end{frame}
\end{document}
